We begin the new Church year in violet vestments, keeping silent the \latin{Gloria}, the angels' hymn of Christmas night. Instead of walking the road to Bethlehem, we stand with Jesus in the Temple courts during Passion Week. He foretells signs in the sun, the moon, and the stars, and nations in distress at the roaring sea. He describes men withering away for fear and the powers of heaven shaken before the Son of Man comes in a cloud with great power and majesty. Hearing those words, we are tempted either to brush them aside or to cling harder to our small earthly security. Yet immediately Jesus turns to His disciples and says, ``look up, and lift up your heads, because your redemption is at hand.''

How can cosmic terror be a command to rejoice, and why does Jesus follow it with the fig tree and all the trees putting forth fruit before summer? Preaching in Rome on this Gospel, St.~Gregory the Great notes that in Scripture the head stands for the mind. To lift up your heads means to gladden your hearts and turn your thoughts toward heaven. Whoever treats this passing world as his permanent home grieves when it shakes, because friendship with the world makes a man an enemy of God. Mortal life is a road, not a home, and he asks how strange it is to grow weary of a hard journey and still wish the road would never end.

On the fig tree putting forth fruit, Gregory reverses what we expect. ``The fruit of the world is its ruin,'' he says. Just as ripening fruit tells a farmer that summer is close, the visible decay of an aging world tells a believer that the eternal summer of God's kingdom is near. Two days before he preached, a sudden whirlwind had struck Rome, uprooting ancient trees, wrecking houses, and throwing down churches, catching men in sudden death while they planned the next morning's business. He asked his people what will happen when the unseen Judge appears in person, if He can shake the foundations of buildings with a single gust of wind. Present calamities are heralds running ahead of the Judge so we may amend our lives before the sentence falls.

Even as this passing age wears out, God does not despise what He made. St.~Cyril of Alexandria explains what Christ means by our redemption being at hand: the resurrection of the dead and the refashioning of creation to what it was at the beginning. When Christ comes openly in majesty, ``this earthly and infirm body shall put off corruption, and shall clothe itself with incorruption by Christ's gift,'' conformed to His glorious body. At the end of today's Gospel, Jesus pledges, ``Heaven and earth shall pass away, but my words shall not pass away.'' As St.~Augustine teaches of the Last Day, the figure of this world passes away, not its nature. Through fire, God renews creation rather than destroying it.

If our redemption is bodily resurrection and a renewed creation, why do our heads stay bent down in dread? The Collect we prayed at the beginning of Mass tells us why. Its opening plea---``Bestir, O Lord, Thy might, we pray Thee and come''---echoes a psalm in which Israel begged God for rescue from invading enemies outside. Here the Church addresses Christ directly and names a closer threat: ``approaching dangers brought on by our sins.'' Sin inside us binds our hearts to a dying world and turns our Redeemer's approach into terror. Unable to free ourselves, we beg Christ to come, that defended by Him we may be rescued and set free by Him we may obtain salvation.

For that reason the Church opens her year with \latin{Ad te levavi animam meam}, ``To thee have I lifted up my soul.'' We sang those words at the entrance and will sing them again at the Offertory. Weighed down by earthly desires, we entrust ourselves to God and hold to the promise from the same psalm that we sang again after the Epistle: ``none of them that wait on thee shall be confounded.'' Twice before the Gospel we also prayed, ``Shew, O Lord, thy ways to me, and teach me thy paths,'' asking for the grace to keep to His road.

St.~Paul puts the same urgency into the Epistle: ``it is now the hour for us to rise from sleep. For now our salvation is nearer than when we believed.'' St.~John Chrysostom sees a quiet drift behind those words. Beginners start the Christian life with eagerness, and as years pass they grow slack. He compares believers to runners in a stadium who push harder once the finish line and the prize come into sight. On Paul's next words---``The night is passed, and the day is at hand''---he pictures neighbors waking one another before dawn at the swallow's cry so sunrise will not catch them in bed.

Paul then changes his wording. He bids us ``cast off the works of darkness,'' and instead of ``works of light,'' he commands us to ``put on the armour of light.'' Ordinary steel armor is heavy and grim, Chrysostom points out, whereas virtue under Christ is light itself, guarding the Christian while making him shine brighter than a sunbeam. Watchfulness becomes a solemn festival, and the believer steps out ``at once soldier and bridegroom.'' Paul names what must be shed before sunrise: rioting and drunkenness, chambering and impurities, contention and envy. Right where today's Epistle ends, he rises from armor to the King Himself: ``put ye on the Lord Jesus Christ.''

A heart caught in old habits cannot wake itself alone. Even while the \latin{Gloria} is silent, we sang Alleluia today: ``Shew us, O Lord, thy mercy; and grant us thy salvation.'' Why ask God to show His mercy, as though it were hidden? St.~Augustine answers that pride tempts us to boast as if our good deeds were our own. God shows His mercy by persuading us that every good we have is His gift. Once a person sees that, pride breaks and he clings to God.

At Communion we will sing from the same psalm that ``the Lord will give goodness: and our earth shall yield her fruit.'' In St.~Augustine's older Latin psalm text, that word was ``sweetness'' where our Missal sings ``goodness,'' and he shows why bare warnings rarely change us. Sin holds us because it still tastes sweet to a divided heart. Before our barren soil can bear fruit, we must stop making excuses and confess our sins truthfully, letting our hearts be turned over by what he calls ``the plough of confession.'' Then God pours a deeper delight in goodness into that opened ground. One who took pleasure in drunkenness comes to rejoice in sobriety, one who delighted in robbing delights in giving, and one who delighted in shows delights in prayer.

God meets us with that cleansing at this altar. In a moment we will stand to profess our faith in the Son born of the Virgin Mary and coming again in glory. Then, over the offerings, the priest will pray silently that these Holy Mysteries, cleansing us by their mighty power, may bring us more pure before God, their Author. After Communion we will ask to receive His mercy in the midst of His temple, ``that with due reverence we may prepare for the coming festival of our redemption.''

Do not wait until the last days before Christmas to turn back to God. Bring into the light of Confession the settled sin you have grown used to excusing, whether excess, impurity, or an old envy and quarrel. Let God's grace break up the hardened ground of your heart and teach you His ways. Advent leads us to Bethlehem, and through Bethlehem to the Day when the Son of Man comes in majesty and our mortal bodies rise incorruptible. Because Christ's words do not pass away, and because He comes to cleanse and feed you at this altar today, you do not have to wither in fear as this world grows old. Wake from sleep, put on the Lord Jesus Christ, and lift up your heads, because your redemption is at hand.
